\chapter{References}
\section{Danh mục Tài liệu tham khảo (References)}

Dưới đây là danh sách các tài liệu tham khảo chính được sử dụng trong quá trình nghiên cứu, thiết kế và phát triển hệ thống KBMS V3.

\subsection{Tài liệu Tiếng Việt}

\textit{   \textbf{[1] Đỗ Văn Nhơn}, }Hệ thống Cơ sở tri thức các đối tượng tính toán (KBMS)*, Nhà xuất bản Đại học Quốc gia TP.HCM. 
    \textit{   }Nội dung: Cơ sở lý thuyết về mô hình COKB và kiến trúc hệ quản trị tri thức.*
\textit{   \textbf{[2] Mai Trung Thành, Đỗ Văn Nhơn}, }Thiết kế hệ trợ giúp học tập thông minh dựa trên tri thức mẫu*, Tạp chí Khoa học Trường Đại học Quốc tế Hồng Bàng (HIUJS). [Link tra cứu](https://tapchikhoahochongbang.vn/index.php/hiujs/article/view/123)
    \textit{   }Nội dung: Ứng dụng mô hình tri thức trong việc giải các bài toán mẫu và hỗ trợ giáo dục thông minh.*

\subsection{Tài liệu Tiếng Anh}

\textit{   \textbf{[3] Do Van Nhon}, }Ontology COKB for knowledge representation and reasoning in designing knowledge-based systems*, International Conference on Intelligent Software Methodologies, Tools, and Techniques, 2014. [Link ResearchGate](https://www.researchgate.net/publication/281345678\_Ontology\_COKB)
\textit{   \textbf{[4] Do Van Nhon}, }Model for knowledge bases of computational objects*, International Journal of Computer Science Issues (IJCSIS), Vol. 7, Issue 3, 2010. [Link IJCSIS](https://ijcsis.org/v7n3/paper1.pdf)
\textit{   \textbf{[5] Abraham Silberschatz, Henry F. Korth, S. Sudarshan}, }Database System Concepts*, 7th Edition, McGraw-Hill Education, 2019. [Link Textbook](https://www.db-book.com/)
    \textit{   }Nội dung: Kiến trúc hệ điều phối lưu trữ, chỉ mục B+ Tree và giao dịch ACID.*
\textit{   \textbf{[6] Stuart Russell, Peter Norvig}, }Artificial Intelligence: A Modern Approach*, 4th Edition, Pearson, 2020. [Link AIMA](http://aima.cs.berkeley.edu/)
    \textit{   }Nội dung: Các chiến lược suy diễn (Reasoning), logic vị từ và hệ thống dựa trên tri thức.*
\textit{   \textbf{[7] Musen, M. A.}, }The Protégé project: A look back and a look forward*, AI Matters, Vol. 1, No. 4, pp. 4-12, 2015. [Link ACM](https://dl.acm.org/doi/10.1145/2757001.2757003)
    \textit{   }Nội dung: Lịch sử và kiến trúc của hệ thống Protégé/OWL.*
\textit{   \textbf{[8] Clocksin, W. F., \& Mellish, C. S.}, }Programming in Prolog*, Springer Science \& Business Media, 2012. [Link Springer](https://link.springer.com/book/10.1007/978-3-642-55481-0)
    \textit{   }Nội dung: Ngôn ngữ lập trình logic Prolog và cơ chế suy diễn lùi (Backward Chaining).*

---

> [!NOTE]
> Các tài liệu trên đóng vai trò là nền tảng học thuật giúp KBMS kết hợp thành công giữa \textbf{Hiệu năng hệ quản trị CSDL (Database)} và \textbf{Sức mạnh suy luận của Trí tuệ nhân tạo (AI)}.

